% Teil II: der durchgerechnete Weg.

\subsection*{Gegeben und gesucht}
Gegeben: Kugel homogen, $I_K=\frac25 m r^2$, Radius $r$, Rollen ohne Gleiten;
Loepingradius $R$ (Bahnradius des Mittelpunkts); Start in Ruhe aus
zusätzlicher Höhe $h$ über dem höchsten Punkt.
Gesucht: $h$, sodass die Kugel im höchsten Punkt gerade nicht herunterfällt.

\subsection*{Ansatz}
Energiesatz von der Ruhe bis zum Loeping-Scheitel (Höhendifferenz $h$);
Kontaktbedingung $N\ge0$ (d.h.\ $v^2\ge gR$) im Scheitelpunkt.

\subsection*{Rechnung}
Energiebilanz (Rollbedingung $\omega=v/r$, $I=\frac25 mr^2$):
\begin{equation}
  m g h = \frac12 m v^2 + \frac12 I \omega^2
  = \frac12 m v^2 + \frac12\,\frac25 m r^2\,\frac{v^2}{r^2}
  = \frac12 m v^2\left(1+\frac25\right)
  = \frac{7}{10} m v^2.
\end{equation}
Damit
\begin{equation}
  v^2 = \frac{10}{7}\,g\,h.
\end{equation}
Kontaktbedingung im Scheitelpunkt ($N=0$ gerade noch,
$\frac{m v^2}{R}=mg$):
\begin{equation}
  v^2 \ge g R
  \;\Rightarrow\;
  \frac{10}{7}\,g\,h \ge g\,R
  \;\Rightarrow\;
  h \ge \frac{7}{10}\,R.
\end{equation}

\begin{ergebnis}
\begin{equation*}
  h_\mathrm{min} = \frac{7}{10}\,R.
\end{equation*}
Die Kugel muss zusätzlich um mindestens \qty{0.70}{} mal den Loopingradius $R$
über dem höchsten Punkt losgelassen werden.
\end{ergebnis}

\begin{kontrolle}
\textbf{Grenzfälle:} Ohne Rotation ($\frac I{mr^2}=0$) wäre $v^2=2gh$ und
$h=\frac12R$. Die Rotation erhöht die nötige Höhe auf $\frac{7}{10}R$ -
plausibel, da ein Teil der Energie "in Rotation steckt". Je kompakter der
Körper (kleineres $I$), desto kleiner das nötige $h$.

\textbf{Einheiten:} $mgh$: $kg\cdot\frac{m}{s^2}\cdot m = kg\,m^2/s^2=J$;
$\frac12 mv^2$:$J$. $h\propto R$ dimensionsmässig konsistent ($m$ kürzt sich). \checkmark

\textbf{Konsistenz mit freiem Fall:} Bei $\Delta h=h$ entsteht $v=\sqrt{\frac{10}{7}gh}$
statt $\sqrt{2gh}$; Faktor $\sqrt{5/7}\approx0.85$ - die Drehenergie bremst die
Zunahme. Vertretbar.
\end{kontrolle}

\begin{merksatz}
Beim Rollen gilt der Energiesatz mit Translations- plus Rotationsenergie;
über die Rollbedingung $\omega=v/r$ wird $v^2(h)$ gewonnen. Zusammen mit der
Kontaktbedingung $v^2\ge gR$ im Scheitel ergibt sich die Mindesthöhe
$h=\frac{7}{10}R$ für die Vollkugel (ohne Rotation: $\frac12 R$).
\end{merksatz}